\section{Exact distance and the closest accurate endpoint}
\label{sec:exact-distance}

A single scalar coordinate describes all the metrics in this section.
Let
\begin{equation}
 b(t)=-\log(1-t^2),\qquad
 \rho(t)=\int_0^t\frac{\sqrt{2(1+u^2)}}{1-u^2}\dd u,
 \qquad -1<t<1.
 \label{eq:rho-def}
\end{equation}
Then $\rho$ is odd, strictly increasing, maps $(-1,1)$ onto $\R$, and
$\rho'(t)^2=b''(t)$. For $0\leq t<1$,
\begin{equation}
 \log\frac1{1-t}\leq\rho(t)\leq\sqrt2\log\frac1{1-t}.
 \label{eq:rho-log}
\end{equation}
Indeed, after multiplying the corresponding derivative inequalities by
$1-t$, these bounds reduce to
$1\leq\sqrt{2(1+t^2)}/(1+t)\leq\sqrt2$.

\subsection{Boxes, rectangular matrix balls, and Jordan intervals}
For a box $(-1,1)^n$ with $F(x)=\sum_i\alpha_i b(x_i)$,
$\alpha_i>0$, integration of the diagonal metric gives
\begin{equation}
 d_F(x,y)^2=\sum_i\alpha_i\bigl(\rho(x_i)-\rho(y_i)\bigr)^2.
 \label{eq:box-distance}
\end{equation}
The transformed straight segment stays in the box and attains this
distance. This scalar product construction is classical; the point of
the next theorem is that leaving a spectral frame cannot shorten the
route from the center.

For $\F\in\{\R,\C\}$ and $p\leq q$, put
\begin{equation}
 \ball_{p,q}^{\F}=\{X\in\F^{p\times q}:\norm{X}_{\op}<1\},\qquad
 \Phi_\alpha(X)=-\alpha\log\det(I_p-XX^*),\quad \alpha>0.
 \label{eq:matrix-ball}
\end{equation}
Write $\sigma^\downarrow(X)$ for its $p$ singular values, in decreasing
order, including zeros.

We also use a Euclidean Jordan algebra $V$, of rank $n$, with unit $e$,
cone of squares $K$, Jordan product $x\circ y$, and trace inner product
$\ip{x}{y}=\tr(x\circ y)$. A Jordan frame $(e_i)_{i=1}^n$ consists of
orthogonal primitive idempotents with sum $e$. The spectral theorem gives
$x=\sum_i\lambda_i(x)e_i$ and $\det x=\prod_i\lambda_i(x)$.
A primitive idempotent is a nonzero element $u$ with $u\circ u=u$
that cannot be written as the sum of two nonzero orthogonal idempotents.
The spectral interval and its barrier are
\begin{equation}
 \mathcal I(V)=\{x:-e\prec_Kx\prec_Ke\},\qquad
 \Phi_\alpha(x)=-\alpha\log\det(e-x^2)
 =\alpha\sum_{i=1}^n b(\lambda_i(x)).
 \label{eq:jordan-interval}
\end{equation}
Here $x^2=x\circ x$ and $u\prec_Kv$ means $v-u\in\interior K$.
The examples include real, complex, and quaternionic Hermitian matrices,
Lorentz algebras, their direct products, and the exceptional Albert
algebra. Rank, rather than the real dimension of $V$, counts the spectral
coordinates.

\begin{theorem}[Exact spectral distance]
\label{thm:exact-distance}
For every $\alpha>0$, the following identities hold:
\begin{align}
 d_{\Phi_\alpha}(0,X)^2
   &=\alpha\sum_{i=1}^p\rho(\sigma_i(X))^2,
   &&X\in\ball_{p,q}^{\F}, \label{eq:matrix-radial-distance}\\
 d_{\Phi_\alpha}(0,x)^2
   &=\alpha\sum_{i=1}^n\rho(\lambda_i(x))^2,
   &&x\in\mathcal I(V). \label{eq:jordan-radial-distance}
\end{align}
More generally,
\begin{align}
 d_{\Phi_\alpha}(X,Y)
 &\geq\sqrt\alpha\,
 \norm{\rho(\sigma^\downarrow(X))-\rho(\sigma^\downarrow(Y))}_2,
 \label{eq:matrix-contraction}\\
 d_{\Phi_\alpha}(x,y)
 &\geq\sqrt\alpha\,
 \norm{\rho(\lambda^\downarrow(x))-\rho(\lambda^\downarrow(y))}_2.
 \label{eq:jordan-contraction}
\end{align}
Here $\rho$ acts coordinatewise and the Jordan eigenvalues are ordered
by signed value. Equality holds in each contraction when the two
endpoints have a common ordered spectral frame. The minimizing path is
linear in the $\rho$ coordinates of that frame.
\end{theorem}
\begin{proof}
First consider matrices. Introduce the Hermitian dilation and a scalar
function
\[
 S(X)=\begin{pmatrix}0&X\\X^*&0\end{pmatrix},\qquad
 g(t)=-\tfrac12\log(1-t^2).
\]
The eigenvalues of $S(X)$ are $\pm\sigma_i(X)$ and $q-p$ zeros, so
$\Phi_\alpha(X)=\alpha\tr g(S(X))$. In an eigenbasis of a Hermitian
matrix $S$ with eigenvalues $\lambda_a$, the spectral-trace Hessian is
\begin{equation}
 D^2\tr g(S)[H,H]
 =\sum_a g''(\lambda_a)|H_{aa}|^2
 +2\sum_{a<b}\frac{g'(\lambda_a)-g'(\lambda_b)}{\lambda_a-\lambda_b}
       |H_{ab}|^2.
 \label{eq:hermitian-Hessian}
\end{equation}
The quotient at equal eigenvalues means $g''(\lambda_a)$.
For completeness, the gradient of $\tr g(S)$ is $g'(S)$.
Differentiating a polynomial in $S$ shows that the $(a,b)$ entry of
$Dg'(S)[H]$ is the displayed divided difference times $H_{ab}$.
The convergent power series of $g$ on $(-1,1)$ justifies the same
calculation here, uniformly on compact spectral subintervals.
Taking its inner product with $H$ gives
\eqref{eq:hermitian-Hessian}, over the real and complex fields alike.
This is the trace-separable case of the spectral Hessian
formula~\cite{LewisSendov2001}.

All divided differences in \eqref{eq:hermitian-Hessian} are positive.
The ordered eigenvalues of an absolutely continuous Hermitian path are
absolutely continuous; almost everywhere their derivatives are the
eigenvalues of the derivative compressed to each repeated eigenspace.
Thus one can choose the eigenbasis within each such space to identify
the diagonal contributions with these derivatives. Applying the formula
to $S(X(t))$ gives
\begin{equation}
 \norm{\dot X(t)}_{\Phi_\alpha,X(t)}^2
 \geq\alpha\sum_i b''(\sigma_i(X(t)))\dot\sigma_i(X(t))^2
 =\alpha\norm{\tfrac{\mathrm d}{\mathrm dt}\rho(\sigma(X(t)))}_2^2
 \quad\text{a.e.}
 \label{eq:singular-speed-contraction}
\end{equation}
The factor is $b''$, because each pair contributes $2g''=b''$.
Singular values are Lipschitz functions of the matrix, and at zero or
repeated singular values the same almost-everywhere rule follows from
the dilation. Integration proves \eqref{eq:matrix-contraction}.
For a fixed singular value decomposition $X=U\diag(\sigma_i)V^*$,
the path
\begin{equation}
 X(t)=U\diag\bigl(\rho^{-1}(t\rho(\sigma_i))\bigr)V^*,
 \qquad 0\leq t\leq1,
 \label{eq:radial-geodesic}
\end{equation}
has constant speed equal to the right side of
\eqref{eq:matrix-radial-distance} to the power $1/2$.
Indeed, restricting $\Phi_\alpha$ to this frame gives
$\alpha\sum_i b(x_i)$. The lower bound therefore is attained.
Interpolation between two coordered transformed spectra gives the other
equality assertion; order is preserved under this interpolation.

Here are the details for Jordan algebras. Relative to a Jordan frame, the
Peirce spaces are $V_{ii}=\R e_i$ and, for $i<j$,
\[
 V_{ij}=\{h:e_i\circ h=e_j\circ h=\tfrac12h,
                  \ e_k\circ h=0\text{ for }k\notin\{i,j\}\}.
\]
At $x=\sum_i\lambda_i e_i$, write the orthogonal Peirce decomposition
of a vector $h$ as $h=\sum_i h_i e_i+\sum_{i<j}h_{ij}$, where
$h_{ij}\in V_{ij}$ and $\norm{e_i}=1$. The Hessian is
\begin{equation}
 D^2\Phi_\alpha(x)[h,h]
 =\alpha\sum_i b''(\lambda_i)h_i^2
 +\alpha\sum_{i<j}
 \frac{b'(\lambda_i)-b'(\lambda_j)}{\lambda_i-\lambda_j}
       \norm{h_{ij}}^2.
 \label{eq:jordan-Hessian}
\end{equation}
One way to verify this formula without restricting the algebra type is
to differentiate the gradient
$\alpha((e-x)^{-1}-(e+x)^{-1})$.
To compute the derivative of inversion at a positive diagonal element
$z=\sum_i z_i e_i$, let $L_z h=z\circ h$ and differentiate
$z\circ z^{-1}=e$. This gives
$L_z D(z^{-1})[h]=-L_{z^{-1}}h$.
The operator $L_z$ is positive definite: it acts by $z_i$ on $V_{ii}$
and by $(z_i+z_j)/2$ on $V_{ij}$.
On the same spaces, $L_{z^{-1}}$ acts by $z_i^{-1}$ and
$(z_i^{-1}+z_j^{-1})/2$, respectively. Solving the differentiated
identity therefore gives the inversion coefficients $-z_i^{-2}$ and
$-(z_i z_j)^{-1}$. The resulting off-diagonal coefficient is
\[
 \frac1{(1-\lambda_i)(1-\lambda_j)}
 +\frac1{(1+\lambda_i)(1+\lambda_j)}
 =\frac{b'(\lambda_i)-b'(\lambda_j)}{\lambda_i-\lambda_j}.
\]
This also gives its continuous value at a repeated eigenvalue.
Every coefficient is positive. Jordan eigenvalues obey the same local
Lipschitz and repeated-eigenspace derivative rule: within the Peirce
algebra for a repeated eigenvalue, diagonalize the corresponding
component of $h$. Equivalently, this follows from the spectral variational
principle applied to $x+th$. Hence, along an absolutely continuous curve,
\[
 \norm{\dot x}_{\Phi_\alpha,x}^2
 \geq\alpha\sum_i b''(\lambda_i(x))\dot\lambda_i(x)^2.
\]
Integrate and interpolate in a fixed Jordan frame exactly as above.
Only the Jordan spectral theorem and Peirce rules have been used, so
the argument includes the exceptional algebra.
\end{proof}

The scale condition in the theorem is only $\alpha>0$, since distance
uses the Hessian. The local-step conclusions need self-concordance.

\begin{proposition}[Self-concordance and parameter]
\label{prop:standard-parameters}
If $\alpha\geq1$, the barriers in \eqref{eq:matrix-ball} and
\eqref{eq:jordan-interval} are standard self-concordant barriers.
Their least gradient parameters are respectively $\alpha p$ and
$\alpha n$. A product with scales $\alpha_j\geq1$ has least gradient
parameter equal to the sum of these parameters.
\end{proposition}
\begin{proof}
The matrix barrier at scale one is the restriction of the positive
Hermitian log-determinant barrier to
$\left(\begin{smallmatrix}I_p&X\\X^*&I_q\end{smallmatrix}\right)$.
The Jordan barrier is the restriction of
$-\log\det u-\log\det v$ to $(u,v)=(e-x,e+x)$.
The log-determinant is self-concordant: after its Hessian isometry at an
interior point, the second and third directional derivatives are
$\sum_i\xi_i^2$ and $-2\sum_i\xi_i^3$, so
$|\sum_i\xi_i^3|\leq(\sum_i\xi_i^2)^{3/2}$.
Affine restriction and addition preserve self-concordance, as does
multiplication by $\alpha\geq1$. Boundary divergence follows from the
vanishing determinant.

The gradient has no component outside the spectral frame, and the
Hessian has no mixed term between that frame and its orthogonal
complement. Thus its squared dual norm is
\begin{equation}
 \norm{\nabla\Phi_\alpha}_{\Phi_\alpha,*}^2
 =\alpha\sum_i\frac{b'(t_i)^2}{b''(t_i)}
 =\alpha\sum_i\frac{2t_i^2}{1+t_i^2}
 <\alpha n_0,
 \label{eq:exact-standard-parameter}
\end{equation}
where $n_0=p$ and $t_i=\sigma_i$ for a matrix ball, and $n_0=n$ and
$t_i=\lambda_i$ for a Jordan interval. In the matrix case this block
structure follows either from \eqref{eq:hermitian-Hessian} or by
differentiating the invariant trace function in singular coordinates.
Letting all $t_i$ tend to $1$ proves that the supremum is $\alpha n_0$.
For products, squared gradient dual norms add, and the limits can be
approached simultaneously.
\end{proof}

\subsection{Distance to an objective target set}
\label{subsec:water-filling}
Consider any finite product of the domains just defined. Each factor
has its own scale $\alpha_j>0$. A nonzero support objective has active
weights $w_i>0$: the positive singular values of each matrix objective,
or the absolute values of the nonzero Jordan eigenvalues of each Jordan
objective. Pool all these weights, indexed by $1\leq i\leq r$, and
attach to each its factor's scale $\alpha_i$. Within a factor all scales
are therefore equal. For a scalar box this simply means
$w_i=|c_i|$ at the nonzero coordinates. Let $W_0=\sum_i w_i$; the
support value is $W_0$ and the analytic center is zero.

\begin{theorem}[The closest accurate endpoint]
\label{thm:water-filling}
For $0<\eps<W_0$, the distance to the entire objective target set is
\begin{equation}
 L_\opt(\eps)^2
 =\min\left\{\sum_{i=1}^r\alpha_i\rho(x_i)^2:
      0\leq x_i<1,\quad\sum_{i=1}^r w_i(1-x_i)\leq\eps\right\}.
 \label{eq:water-filling}
\end{equation}
The scalar allocation problem has a unique minimizer, all coordinates
lie in $(0,1)$, and its constraint is active. There is a unique
$\lambda>0$ such that
\begin{equation}
 2\alpha_i\rho(x_i)\rho'(x_i)=\lambda w_i\quad(1\leq i\leq r),
 \qquad \sum_iw_i(1-x_i)=\eps.
 \label{eq:water-KKT}
\end{equation}
An aligned endpoint with these coordinates and all inactive coordinates
zero attains the minimum distance in the original domain. The path
linear in its transformed spectral coordinates is globally shortest
to the target set.
\end{theorem}
\begin{proof}
Distances in a product metric add in squares. To see this directly, the
length of a product curve is the integral of the Euclidean norm of its
component speeds, which is at least the Euclidean norm of the component
lengths. Constant-speed shortest component curves attain equality; such
curves from zero were constructed in \cref{thm:exact-distance}.

For a matrix factor, von Neumann's trace inequality gives
$\ip{C}{X}\leq\sum_iw_i\sigma_i(X)$, with weights and singular values
ordered. Thus replacing an endpoint by one in the objective's singular
frame can only improve its objective and leaves its distance unchanged.
Inactive singular values can then be set to zero, reducing the distance.
For a Jordan factor, if $(e_i)$ and $(f_j)$ are frames of $c$ and $x$,
the entries $m_{ij}=\ip{e_i}{f_j}$ are nonnegative and their row and
column sums are one. Therefore
\[
 \ip{c}{x}=\sum_{i,j}\mu_i\lambda_j m_{ij}
 \leq\sum_i|\mu|_i^\downarrow|\lambda|_i^\downarrow.
\]
The last inequality is the rearrangement inequality for a doubly
stochastic matrix. Assigning magnitudes to the objective frame in this
order and using the signs of $\mu_i$ attains the bound, at the same
radial distance. Zero-weight coordinates can again be removed.
Conversely every scalar vector in \eqref{eq:water-filling} gives an
aligned feasible endpoint with precisely that objective error and
radial distance. This proves the reduction.

For $x\geq0$, $\rho'>0$, $\rho''\geq0$, and
$(\rho^2)''=2((\rho')^2+\rho\rho'')>0$.
The objective is therefore strictly convex. It tends to infinity if
any coordinate tends to $1$. A nonempty bounded sublevel set is compact
in $[0,1)^r$, so a minimizer exists and is unique.
Its accuracy constraint must be active: otherwise reduce one positive
coordinate slightly. Slater's condition holds by taking every
coordinate sufficiently close to $1$. The KKT multiplier is positive,
since a zero multiplier would force every coordinate to be zero, which
violates $\eps<W_0$. At $x_i=0$ the objective's derivative is zero,
whereas the accuracy constraint contributes $-\lambda w_i<0$;
stationarity with a nonnegative multiplier for $-x_i\leq0$ is then
impossible. Every coordinate is positive, and stationarity gives
\eqref{eq:water-KKT}.

The function $(\rho^2)'$ increases continuously from zero to infinity.
For a fixed $\lambda>0$, each equation in \eqref{eq:water-KKT} thus has
one solution. Their weighted residual decreases continuously and
strictly from $W_0$ to zero as $\lambda$ increases from zero to infinity.
This proves uniqueness of $\lambda$ and supplies a scalar search.
The minimizing path follows from \cref{thm:exact-distance} and the
product rule. Within each factor the minimizer is ordered with its
weights because its scales agree, so the common-frame construction
applies without an ordering issue.
\end{proof}

The uniqueness assertion concerns the scalar allocation. It is not
needed to assert uniqueness of all ambient matrix representations at
repeated objective eigenvalues. If $\eps\geq W_0$, the analytic center
has zero cost. For $\eps=0$, no strictly interior point is accurate, and
$L_\opt(\eps)\to\infty$ as $\eps\downarrow0$.

\begin{corollary}[Logarithmic allocation and bounded moves]
\label{cor:water-movement}
Let $0<\eps<W_0$. Define
\begin{equation}
 D_{\log}(\eps)^2=
 \min\left\{\sum_i\alpha_i s_i^2:s_i\geq0,
                    \quad\sum_i w_i e^{-s_i}\leq\eps\right\}.
 \label{eq:log-water-filling}
\end{equation}
Then
\begin{equation}
 D_{\log}(\eps)\leq L_\opt(\eps)\leq\sqrt2D_{\log}(\eps).
 \label{eq:log-comparison}
\end{equation}
Its unique optimizer satisfies
\begin{equation}
 s_i=\mathcal W(\gamma w_i/\alpha_i),\qquad
 \sum_i\alpha_i\mathcal W(\gamma w_i/\alpha_i)=\gamma\eps,
 \quad\gamma>0,
 \label{eq:Lambert}
\end{equation}
where $\mathcal W$ is the nonnegative inverse of $s\mapsto s e^s$.
For $0<R<1$, if all scales are at least one, the minimum number
$T_R^\star(\eps)$
of feasible forward-$R$ moves from zero to any accurate endpoint obeys
\begin{equation}
 \frac{L_\opt(\eps)}{-\log(1-R)}
 \leq T_R^\star(\eps)
 \leq\left\lceil\frac{L_\opt(\eps)}{\log(1+R)}\right\rceil.
 \label{eq:exact-bounded-movement}
\end{equation}
\end{corollary}
\begin{proof}
Set $s_i=-\log(1-x_i)$ and apply \eqref{eq:rho-log} to every feasible
allocation. Strict convexity and the KKT conditions for
\eqref{eq:log-water-filling} give
$\alpha_i s_i=\gamma w_i e^{-s_i}$ with all $s_i>0$.
Inverting and summing these identities gives \eqref{eq:Lambert}.
Uniqueness follows, for example, by observing that
$\sum_iw_i\exp[-\mathcal W(\gamma w_i/\alpha_i)]$ decreases strictly
from $W_0$ to zero; the zero root of the multiplied equation in
\eqref{eq:Lambert} is excluded by $\gamma>0$.
Finally, \cref{prop:standard-parameters,cor:optimal-movement} apply,
and the shortest curve was constructed in \cref{thm:water-filling}.
\end{proof}

The construction is explicit once the objective frame is available.
Invert $(\rho^2)'$ coordinatewise and use doubling and bisection for the
positive multiplier. An upper bracket whose residual is at most the
requested tolerance gives an accurate endpoint. Along the resulting
path, checkpoints at arc intervals $\log(1+R)$ satisfy the forward
local-norm restriction. Finite-precision evaluation must preserve these
two inequalities, for instance by using strict accuracy and step-radius
margins. We make no arithmetic-complexity claim for this scalar search
or for obtaining the spectral frame.
